\documentclass[12pt]{article}
\usepackage{amsmath, amssymb}

\title{CSE 400 -- Lecture 21 \& 22 Scribe\\
Inequalities and Tail Bounds}
\date{}

\begin{document}

\maketitle

\section{Introduction to Tails}

\textbf{Definition (Tail of a Random Variable)}\\
The tail of a random variable $X$ is:
\[
P(X > x)
\]

\textbf{Why do we care about tails?}
\begin{itemize}
    \item Jobs delayed more than 24 hours
    \item Packets dropped due to full buffer
\end{itemize}

\textbf{Key Idea}
\begin{itemize}
    \item Mean $\rightarrow$ average behavior (easy)
    \item Tail $\rightarrow$ rare extreme events (hard)
\end{itemize}

\textbf{Why are tails hard to compute?}\\
Requires summing many small probabilities

\section{Tail Examples}

\subsection{Example 1: Jobs distributed across servers}

\textbf{Question}\\
Suppose you’re distributing $n$ jobs among $n$ servers at random.\\
What is the probability that a particular server gets $\geq k$ jobs?

\textbf{Intuition}
\begin{itemize}
    \item Jobs are assigned randomly
    \item Most servers get few jobs
    \item Occasionally, one server gets many jobs
\end{itemize}

$\rightarrow$ This is a tail event (rare overload)

\textbf{Model}
\[
X \sim \text{Binomial}(n, p)
\]

\textbf{Exact Probability}
\[
P(X \geq k) = \sum_{i=k}^{n} \binom{n}{i} p^i (1-p)^{n-i}
\]

\subsection{Example 2: Poisson arrivals}

\textbf{Question}\\
Jobs arrive to a datacenter according to a Poisson process with rate $\lambda$ jobs/hour.\\
What is the probability that $\geq k$ jobs arrive during the first hour?

\textbf{Intuition}
\begin{itemize}
    \item Jobs arrive randomly over time
    \item Most hours have moderate arrivals
    \item Sometimes, we see a burst of arrivals
\end{itemize}

$\rightarrow$ This is a tail event (sudden spike in load)

\textbf{Model}
\[
X \sim \text{Poisson}(\lambda)
\]

\textbf{Exact Probability}
\[
P(X \geq k) = \sum_{i=k}^{\infty} e^{-\lambda} \frac{\lambda^i}{i!}
\]

\section{Why Tail Bounds?}

\textbf{Observation}\\
Exact tail probabilities are often hard to compute

\textbf{Exact Computation}
\[
P(X \geq k) = \sum_{i=k}^{\infty} e^{-\lambda} \frac{\lambda^i}{i!}
\]
\[
P(X \geq k) = \sum_{i=k}^{n} \binom{n}{i} p^i (1-p)^{n-i}
\]

$\rightarrow$ These are complicated, long sums

\textbf{Idea}
\begin{itemize}
    \item Find an upper bound
    \item Easier to compute
    \item Still gives a guarantee
\end{itemize}

\textbf{Goal}
\[
\text{Exact: } P(X \geq k) = ?
\]
\[
\text{Bound: } P(X \geq k) \leq \text{something simple}
\]

\textbf{Key Idea}\\
Approximate safely instead of computing exactly

\section{Running Example}

We will:
\begin{itemize}
    \item Develop progressively better (tighter) tail bounds
    \item Test all bounds on the same example
\end{itemize}

\textbf{Experiment}\\
Flip a fair coin $n$ times

\textbf{Random Variable}
\[
X = \text{number of heads}
\]
\[
X \sim \text{Binomial}\left(n, \frac{1}{2}\right)
\]

\textbf{Mean}
\[
E[X] = \frac{n}{2}
\]

\textbf{Question}
\[
P\left(X \geq \frac{3}{4}n\right)
\]

$\rightarrow$ Tail region: unusually many heads

\section{Markov’s Inequality}

\textbf{Key Idea}
\begin{itemize}
    \item Uses only the mean
    \item Large values are rare
\end{itemize}

\textbf{Example Idea}\\
Average = 10 $\rightarrow$ few values can be 100+

\textbf{Markov Inequality (Statement)}\\
For a non-negative random variable $X$:
\[
P(X \geq a) \leq \frac{E[X]}{a}
\]

\textbf{Interpretation}
\begin{itemize}
    \item Probability that $X$ is large
    \item Controlled only by its mean
\end{itemize}

\textbf{Proof (Step-by-step)}

We start with:
\[
E[X] = \sum_x x \cdot P(X = x)
\]

Split into two parts:
\[
x < a \quad \text{and} \quad x \geq a
\]

\[
E[X] \geq \sum_{x \geq a} x \cdot P(X = x)
\]

Since $x \geq a$:
\[
E[X] \geq \sum_{x \geq a} a \cdot P(X = x)
\]

\[
E[X] \geq a \sum_{x \geq a} P(X = x)
\]

\[
E[X] \geq a \cdot P(X \geq a)
\]

Divide both sides by $a$:
\[
P(X \geq a) \leq \frac{E[X]}{a}
\]

\section{Application to Running Example}

\[
X \sim \text{Binomial}\left(n, \frac{1}{2}\right), \quad E[X] = \frac{n}{2}
\]

We want:
\[
P\left(X \geq \frac{3}{4}n\right)
\]

Apply Markov:
\[
P(X \geq a) \leq \frac{E[X]}{a}
\]

Substitute:
\[
P\left(X \geq \frac{3}{4}n\right) \leq \frac{n/2}{3n/4}
\]

\[
= \frac{1}{2} \cdot \frac{4}{3} = \frac{2}{3}
\]

\textbf{Result}
\[
P\left(X \geq \frac{3}{4}n\right) \leq \frac{2}{3}
\]

\textbf{Observation}
\begin{itemize}
    \item Very loose bound
    \item Only uses mean
    \item Does not use variance or distribution shape
\end{itemize}

\section*{Next Topics}

\begin{itemize}
    \item Chebyshev’s Inequality
    \item Chernoff Bounds (Poisson tail)
    \item Jensen’s Inequality
    \item Convex functions
\end{itemize}

\end{document}